\section{Localization and finite-level periods}\label{sec:localization}

We now connect the local pairing of Section~\ref{sec:local} with a
Whittaker coefficient of global localization.  The comparison has two
steps.  This section expresses that coefficient as an exponential period
on a stack of unipotent bundles.  Sections~\ref{sec:unipotent}
and~\ref{sec:ran} identify the period with semi-infinite cohomology after
integration over additional markings.  The integration is essential:
semi-infinite cohomology at one fixed finite set is not itself the global
period.

Continue with the notation $G,A,s,C_T,d_T$ of
Sections~\ref{sec:loop}--\ref{sec:local}.  Set
\[
 \delta_N=-\chi(X,\mathfrak n_0),\qquad
 l_N=\det R\Gamma(X,\mathfrak n_0).
\]
We compose determinant lines and their dimension shifts using graded
multiplicativity, including the relative-lattice identifications in
\eqref{eq:L6}.  For a tuple $T$ over a smooth parameter scheme $S$, write
\[
 \operatorname{Loc}_T:C_T\longrightarrow
       D(L_G^s\boxtimes\mathcal O_S),\qquad
 \Lambda_T=\operatorname{Loc}_T\otimes l_N^{-1}[\delta_N].
\]
The target is the category on $Y_G\times S$; the parameter is retained.
We construct $\operatorname{Loc}_T$ below, specifying its left-module
normalization.  The functor $\operatorname{oblv}_W:Q_S\to D(S)$ is the
common underlying-complex functor of \eqref{eq:walg-forgetful} and
Theorem~\ref{walg:common-target};
its semi-infinite computation uses \eqref{eq:L6} with the same
left-crystal normalization.  The coefficient $R_T$ retains the marked
parameters, and the pairing is that of \eqref{eq:W2}.

Write $\int_{V\supseteq T}$ for relative Ran integration with dualizing
normalization.  Its charts are all tuples enlarging $T$, with their
incidence maps, including coincidences.  It can be computed by direct
images from these tuple and incidence schemes followed by the colimit
of their $!$ systems; equivalently, one uses the colimit presentation
with closed direct images.  These presentations are pseudo-proper,
so the direct images in this calculation agree for $*$ and $!$.
Consequently the integral is continuous, commutes with $!$ base change,
and satisfies Fubini for two successive enlargements.

\begin{theorem}[Period comparison]\label{loc:period-comparison}
Fix $G$ and its level as in Section~\ref{sec:loop}, and let $T$ be a
nonempty tuple over a smooth parameter scheme $S$.  There is a natural
isomorphism of continuous $D(S)$-bilinear pairings
\begin{equation}\label{eq:T1}
 \left\langle R_T\Lambda_T(M),W'\right\rangle_S
 \simeq
 \int_{V\supseteq T}
 \operatorname{oblv}_W\!\left(
 F_V(\operatorname{ins}_{T,V}M,
             \operatorname{ins}_{T,V}W')\right),
 \qquad M\in C_T,\quad W'\in d_T.
\end{equation}
It respects the ordinary and crystalline parameter comparisons,
including permutations, repetitions, and diagonal mergers of marks.
\end{theorem}

The proof occupies the remainder of this section and
Sections~\ref{sec:unipotent}--\ref{sec:ran}.

We will also use the universal homological contractibility of the Ran
space of a connected curve \cite{GaitsgoryContractibility}, in its
relative form with a prescribed tuple contained in the set.  Indeed,
union with that tuple maps ordinary Ran to the containing-set Ran;
forgetting the containment and then taking union is the identity on the
latter.  Its augmentation is thus a retract of the ordinary Ran
augmentation, also after base change.  The augmentation therefore
integrates a $!$-pullback from the base to that same object.  Several
prescribed tuples are handled by their union.  Our prescribed tuples
are nonempty.

\subsection{Level structures with moving markings}

Fix a smooth affine chart $S$ of a tuple space.  Let $D_{\mathrm{lv}}$
be either zero or an effective thickening of the entire tuple $T$;
positive multiples of the sum of its graphs form sufficient choices.
Consider the smooth group schemes over $X\times S$
\[
 \mathfrak G=(G^{P_0})^{D_{\mathrm{lv}}},\qquad
 \mathfrak N=N_0^{D_{\mathrm{lv}}}.
\]
The superscript means dilation imposing the identity along the level
divisor.  For $N_0$, root coordinates are the corresponding root lines
twisted by $-D_{\mathrm{lv}}$; their multiplication is defined because
products have at least the required vanishing.  For $G^{P_0}$,
smoothness follows in \'etale coordinates about the identity by
dividing the coordinates by a local equation of the divisor.  Away from
that divisor the group scheme is unchanged.

Let $B_G^{\mathrm{lv}}$ and $B_N^{\mathrm{lv}}$ be their bundle stacks.
They classify bundles with the indicated level framing in the inner
presentation.  In particular $B_G^0=Y_G\times S$, by contraction with
$P_0$.  Curve deformation theory makes both stacks smooth over $S$.
Write $\mathfrak g^{\mathrm{lv}}$ and
$\mathfrak n^{\mathrm{lv}}=\mathfrak n_0(-D_{\mathrm{lv}})$ for the Lie
bundles.  The twist on $B_G^{\mathrm{lv}}$ is pulled back from $Y_G$.
For the natural map
\[
 p:B_N^{\mathrm{lv}}\longrightarrow B_G^{\mathrm{lv}},
\]
its pullback is canonically trivial: the root filtration with fixed
Cartan bundle $T_0$ identifies its determinant ratio with the ratio at
$P_0$.

For a larger tuple $E\supseteq T$, let
$H_{G,E}^{\mathrm{lv}}$ and $H_{N,E}^{\mathrm{lv}}$ be the regular arc
groups of these level group schemes.  The meromorphic Lie algebras do
not change, since the level divisor is supported on $T\subseteq E$.
We use the Harish--Chandra module categories of
Proposition~\ref{loop:hc-renormalization}; on the $N$ side no character condition is imposed here.
The affine extension on the $G$ side has central scalar $-s$ and splits
canonically on $N$.

\begin{lemma}\label{loc:levels}
The level Harish--Chandra categories have the compact presentations,
continuous functors, and ordinary and crystalline parameter comparisons
of Propositions~\ref{loop:hc-renormalization} and~\ref{loop:crystal}.  Restriction between two nonzero levels supported on
the whole old tuple is fully faithful when compared through a common
deeper level.  It preserves compact objects.  These comparisons remain
valid in the presence of extra markings, including collisions with the
old ones.
\end{lemma}

\begin{proof}
First take the underlying quasi-coherent calculation, denoted by $b$ in
Section~\ref{sec:loop}.  The level subgroup is the kernel of evaluation
from regular arcs to the level divisor.  At sufficiently deep finite
jet stages this evaluation is a smooth quotient.  Its differential is
surjective by restriction between finite divisors, whose modules are
finite locally free.  Surjectivity can be checked on geometric points:
a jet of a smooth group extends over local Artin thickenings; at an
extra distinct point choose the identity.  The resolution-property
argument in Proposition~\ref{loop:hc-renormalization} applies to this kernel.  Namely, the ambient
jet group has an embedding with quasi-affine homogeneous space, and
the additional quotient imposed here is affine.  The compact
presentation and bounded-object arguments therefore carry over.

One must distinguish two uses of unipotence.  At an added point away
from $T$, the level subgroup includes full arcs and need not itself be
pro-unipotent.  The \emph{difference between deeper levels on the old
tuple}, however, is smooth unipotent.  In the coordinates
$u=P(z)$ of Section~\ref{sec:loop}, use positive powers of $u$; successive
kernels are additive groups from square-zero thickenings.  Thus
forgetting between these levels forgets strong equivariance for a
unipotent quotient.  Normalized smooth pull and affine-space de Rham
acyclicity give full faithfulness, and the finite relative Chevalley
resolution gives preservation of compactness.  Full-arc factors at the
added points do not affect this argument.

For crystal comparisons, take two lifts of the parameter tuple across
a nilpotent ideal.  Their full formal disks and meromorphic objects,
including the affine data of \eqref{eq:L1}, agree.  Pass to common
thickenings of their level divisors.  Such thickenings are cofinal by
nilpotence, and the relevant divisor quotients are finite locally free;
locally this follows by division of their monic equations.  Differences
on the old levels again have unipotent quotients, since their defining
ideals in the larger Artin jet rings are nilpotent.

To include the full-arc factors, first divide by a common deeper normal
congruence on the entire tuple $E$.  The finite relative D-module
equivariant diagrams now compare by their relative de Rham groupoids,
together with the formal Lie transport in \eqref{eq:L1}.  In particular,
this comparison uses the specified identification of the actions; an
isomorphism of reduced spaces alone would not specify a comparison of
Lie modules.  Restriction to the common level realizes the same
comparison.  The finite equivariance calculation of
Proposition~\ref{loop:strong-action}, or
the relative Chevalley resolution for a unipotent difference, identifies
the corresponding module functors.  Zero level compares directly by
the full-disk identification.  Refinement to a third common level
proves the cocycle identities, including on multiple infinitesimal
diagonals.  This gives precisely the crystal parameter operations.
\end{proof}

\subsection{The localization algebroid}

We next construct localization on smooth bundle charts.  This also
identifies the maps that will be needed when a Grassmannian kernel
moves in its parameters.

Let $U$ be a smooth scheme over the parameter base with a map to one
of the bundle stacks and with full disk framings along $E$, compatible
with its level.  Write $\underline\ell_U$ for the meromorphic Lie
algebra, with its affine extension on the $G$ side.  Infinitesimal
replacement of the framed gluing data gives the transitive algebroid
\[
 \mathcal T_U
   =T_{U/\mathrm{base}}
       \mathop{\times}_{T_{\mathrm{Framed}/\mathrm{base}}}
                   \underline\ell_U.
\]
An element is a pair $(v,x)$ for which the replacement by $x$ induces
the deformation $v$ of the framed bundle.  The bracket is the bracket
of the right action, including differentiation of Lie sections.
Its anchor kernel is the algebra of outside currents
\[
 \mathfrak g^{\mathrm{lv}}_{P,\mathrm{out}}
       =\Gamma(X-E,\mathfrak g^{\mathrm{lv}}_P),
\]
together with the center when present; the analogous assertion holds
for $N$.

Here is why this remains an ordinary algebroid calculation despite the
Tate notation.  Every infinitesimal framed deformation can be
trivialized away from $E$: that complement is affine, so its relative
$H^1$ vanishes locally on the chart.  Such a trivialization supplies
the Laurent replacement and proves transitivity.  Its ambiguity is
exactly an outside current.  The latter is a union of locally free
finite-pole pieces with locally free successive quotients, after
starting at a sufficiently large pole bound with vanishing $H^1$.
On an affine chart, split the successive quotients: the outside
currents are then a countable direct sum of finite projective modules.
The projective form of PBW therefore applies.  Full framings can be chosen by first choosing finite jets and then lifting
through smooth jet projections with unipotent kernels on affine test
schemes.  An automorphism preserving a full framing is trivial.

The determinant line splits the affine extension on outside currents:
an outside replacement identifies the bundle back with itself, and
there is a unique lift acting trivially on its determinant.  The
relative determinant of lattices gives this action.  Differentiation
on the determinant torsor gives a map from $\mathcal T_U$ to the
relative twisted differential operators $\mathscr D_U^s$, taking the
center to $-s$.

\begin{proposition}\label{loc:construction}
Localization with level is the continuous functor whose calculation
on such a chart is
\begin{equation}\label{eq:T2}
 \operatorname{Loc}_{G,E}^{\mathrm{lv}}(M)|_U
   =\mathscr D_U^s\otimes^L_{U(\mathcal T_U)}M_U,
 \qquad (v,x)m=v(m)+xm.
\end{equation}
Here $M_U$ is the module pulled to the framed chart and the enveloping
algebra uses the specified central quotient.  Omitting the center and
twist gives $\operatorname{Loc}_{N,E}^{\mathrm{lv}}$.
These functors commute with the parameter operations of
Proposition~\ref{loop:crystal}, with pullback between smooth bundle charts, and with
the deeper-level comparisons of Lemma~\ref{loc:levels}.
At zero level divisor, $\operatorname{Loc}_{G,T}^0$ is
$\operatorname{Loc}_T$.
\end{proposition}

\begin{proof}
Formula~\eqref{eq:T2} is a formula for left modules in the
angle-bracket convention.  Its underlying $\mathcal O$ calculation is
the Chevalley chain complex of the outside Lie algebra with
coefficients in $M_U$, without an additional shift inside the
angle brackets.  Indeed, as a right $U(\mathcal T_U)$-module the TDO is
the quotient by the outside kernel, using its determinant splitting.
Resolve that kernel and apply PBW; equivalently use the Spencer
resolution and the anchor exact sequence.  Projective outside
algebras of this form cause no completion: the associated graded Koszul
calculation is the direct union of the finite calculations.

For a map of smooth charts $f:U'\to U$, not necessarily a smooth map,
the algebroid on $U'$ is the inverse-image algebroid, formed by taking
the fiber product with $T_{U'/\mathrm{base}}$.  The ordinary TDO
transfer defines the pullback comparison for \eqref{eq:T2}.  On its
generating section it sends module generators to their pulls; the
relations for $(v,x)$ become differentiation and the inverse-image
algebroid action.  Derived tensor or Spencer resolutions give the
map on complexes.  On forgetting differential operators it is exactly
derived pullback of the outside chain complex, hence an isomorphism.
Composition of transfers proves coherence.

Change of a disk frame by an arc element acts on $M_U$ through its
Harish--Chandra equivariance.  On the algebroid it acts by conjugation
and the logarithmic derivative.  The strong differential equality,
with the affine splitting on frames, identifies \eqref{eq:T2} in the
two frames.  This descends the chart formulas.  The same chart
calculation proves compatibility with forgetting additional level,
where the bundle-side operation is $!$-pullback.

We give the parameter verification since it will be used for arbitrary
kernels.  Work with relative differential operators on smooth bundle
charts.  At zero level, the full frame prestack and the Lie formula
identify across infinitesimal lifts of the marked tuple.  At nonzero
level, pull both formulas to a common deeper level as in
Lemma~\ref{loc:levels}.  The two forgetful maps agree on reduction and
supply the same crystalline comparisons; their normalized pulls are
fully faithful because the level differences are unipotent torsors.
The chart formula already agrees on those pulls.  Descent and further
common refinements give all multiple-diagonal cocycles.  For twisted
differential operators, a change of determinant frame by a nil-unit
uses Kummer equivariance, that is, its logarithmic derivative with
scalar $-s$.

Ordinary changes of tuple schemes are checked by derived base change
in the same formulas.  On the inducing compact presentations, PBW and
Chevalley terms are locally flat.  Base changes between smooth schemes
have finite Tor dimension; the cofinal infinitesimal-diagonal
projections used for the crystal test are flat.  Thus these resolutions
compute the asserted comparisons.  This argument works on all smooth
charts of a bundle-stack exhaustion and their common-level pulls; it
does not require a tensor-product assertion for the limit over
non-quasi-compact charts.

Finally take the continuous extension from the compact presentations
of Proposition~\ref{loop:hc-renormalization}.  The formulas already use colimits and derived
tensors.  Their agreement with bounded-module calculations follows by
filtered colimits in bounded ranges and the displayed resolutions.
At zero level the construction is the usual Beilinson--Drinfeld
localization: twist on the frame torsor, descend by positive-loop
splitting, and induce along the uniformization algebroid.  Pullback to
a chart gives \eqref{eq:T2}, or equivalently the outside coinvariants.
This also fixes its normalization, including the ratio rigidification
at $P_0$.  The inner presentation changes only the chosen frames.
\end{proof}

The arbitrary-level localization functor and its crystalline
transformation are constructed in \cite[Section~12.1, especially
Equation~(12.4)]{GLCII}.  Its comparison with the quasi-coherent
construction after insertion and integration is
\cite[Theorems~12.1.5 and~12.1.8]{GLCII}.  We use the standard
localization construction here.  The moving-parameter and
arbitrary-kernel comparisons required below are proved by the
algebroid calculation, and the period identity \eqref{eq:T1} will be
proved in the following sections.

\subsection{Propagation and restriction}

Adding a marked point enlarges both the outside algebra and the
principal parts.  Their quotient records the same global cohomology;
this is the reason localization is unchanged by insertion.

Fix $T\subseteq U'\subseteq E$ and the level divisor supported on $T$.
Insertion in the level categories is denoted by
$\operatorname{ins}_G$ or $\operatorname{ins}_N$.  At $b$ it is
\eqref{eq:L9} with the arcs of the level group scheme.  More explicitly,
use the intermediate Laurent algebra completed along $E$ but allowing
poles only on $U'$, restrict the old module to it, and then induce to
the Laurent algebra at $E$.  The cocycle identifies under this map by
the relative-residue formula, equivalently by the relative determinant.
For each pole bound, principal parts of the intermediate algebra
modulo its regular level lattice are the same finite
divisor-supported module as at $U'$.  This statement uses the sum of
graphs and remains true at collisions.

The induction bar with the common arc subgroup gives the bounded
module formula and its maps.  Its pole quotient is flat by the jet
localization calculation following \eqref{eq:L9}, since the level Lie
bundle is locally free on the curve.  The formula on compact
presentations therefore extends continuously and is transitive.
Arc representations are pulled by the map from arcs at $E$ to those
at $U'$.  For an infinitesimal parameter comparison, the intermediate
algebra compares by inverting the old divisor equation.  Deepening a
level on $T$ leaves this Lie induction unchanged and only restricts
its Harish--Chandra equivariance.  Thus insertion has all the parameter
and level comparisons just constructed.

\begin{proposition}\label{loc:propagation}
With parameters retained, there are natural isomorphisms
\begin{equation}\label{eq:T3}
 \operatorname{Loc}_{G,U'}^{\mathrm{lv}}(M)
 \simeq
 \operatorname{Loc}_{G,E}^{\mathrm{lv}}(
                      \operatorname{ins}_G M),
\end{equation}
and the analogous isomorphisms for $N$.  They respect transitivity,
ordinary and crystalline parameter changes, and deeper levels.
\end{proposition}

\begin{proof}
On an $E$-framed chart, restriction of an outside trivialization from
$X-U'$ to $X-E$ maps the old localization algebroid to the new one.
Its Laurent vector lies in the intermediate algebra: in vector
representations, clear powers of the old divisor equation before
completing along $E$.  This also proves the assertion when points
collide.  The central splittings and the maps to TDOs agree, since
both use the determinant of the same replacement.  Sending module
generators into the induction therefore defines the comparison.

Test the map on the inducing generators of
Proposition~\ref{loop:hc-renormalization}.  Filter the
outside-chain calculation by total PBW degree, counting the number
of chain factors as well.  For an inducing representation $W$, the
associated graded is $W$ pulled to the chart, tensored with
\[
 \operatorname{Sym}\!\left(
  \operatorname{Cone}\bigl(
   \mathfrak g^{\mathrm{lv}}_{P,\mathrm{out}}
     \longrightarrow\operatorname{PP}(
                              \mathfrak g^{\mathrm{lv}}_P)
                     \bigr)\right).
\]
The outside-current term in this cone is in shift $[1]$.
The meromorphic resolution identifies the cone with
$R\Gamma(X,\mathfrak g^{\mathrm{lv}}_P)[1]$ for either tuple,
and the comparison induces the identity on this object.  The
exhaustive increasing PBW filtration gives the isomorphism.  The
argument for $N$ is identical, and all its maps are the parameter
maps of Proposition~\ref{loc:construction}.
\end{proof}

Restriction of Lie modules from $G$ to $N$, denoted
$\operatorname{Res}$, is ordinary bounded restriction on compact
inputs at $b$, with values in the renormalized $N$ category by bounded
comparison, and is then extended continuously.  One may construct it
by the map of Harish--Chandra induction monads: restrict the group
representation and use the enveloping-algebra map.  On finite
representation generators this is the ordinary restriction map;
PBW commutes with filtered colimits in bounded ranges.  The
forgetful functors of Proposition~\ref{loop:hc-renormalization}
show that it is $t$-exact at $b$.
These descriptions also prove the common-level crystal comparisons.

We record more precisely the filtration used later to control this
restriction.  For a $G$-inducing generator with representation $W$,
count only the PBW factors outside the $N$ algebra.  Its $i$th graded
piece, as an $N$ module, is
\begin{equation}\label{eq:T4}
 \operatorname{Ind}_{\operatorname{Lie}H_N^{\mathrm{lv}}}
                   ^{\ell_N^{\mathrm{loop}}}
 \left(W\otimes
  \operatorname{Sym}^i\operatorname{PP}(
                     \mathfrak g^{\mathrm{lv}}/
                     \mathfrak n^{\mathrm{lv}})\right).
\end{equation}
Here $\ell_N^{\mathrm{loop}}$ has no center and $W$ is restricted to
$H_N^{\mathrm{lv}}$.  Order $N$ operators first in PBW.  Commuting an
arc $N$ vector past an outside operator gives its adjoint action
modulo $N$ and the regular lattice; all other terms have smaller
outside count.  Commutators of outside operators, including central
terms, also lower that count.  This proves the intrinsic filtration.
The inducing representations in \eqref{eq:T4} are exhausted by
locally free pieces at finite pole and representation bounds:
finite principal-part bounds are arc-stable, and one takes their
symmetric and tensor powers with $W$.  Consequently the filtration
is valid in the renormalized comparison and commutes with scalar
extension and the crystal maps.

The enveloping formula gives a natural map
\[
 \operatorname{ins}_N\operatorname{Res}(M)
       \longrightarrow
       \operatorname{Res}(\operatorname{ins}_G M).
\]
The map of localization algebroids on $B_N^{\mathrm{lv}}$, together
with the determinant trivialization there, gives
\begin{equation}\label{eq:T5}
 \operatorname{Loc}_{N,E}^{\mathrm{lv}}(
                              \operatorname{Res}M)
       \longrightarrow
 p^!\operatorname{Loc}_{G,E}^{\mathrm{lv}}(M).
\end{equation}
These maps commute with propagation \eqref{eq:T3}, with the displayed
insertion map for restriction, and with the ordinary, crystalline,
and deeper-level comparisons.  They are maps induced by the same
derived tensor construction \eqref{eq:T2}.

\subsection{Localization of an arbitrary Grassmannian kernel}

Let $M\in C_T$ and let
$K\in D_{\mathcal L^{-s}}(\operatorname{Gr}_T)$ be supported on a
finite closed bound $Z$, which may be taken proper over $S$.  Choose a
level divisor on $T$ sufficiently deep that its congruence subgroup
acts trivially on $Z$ and that all determinant data factor through
the resulting finite jets.  Such a choice exists: for Laurent elements
in a fixed bound, conjugation carries sufficiently deep congruences
into positive arcs.  Then $K\star M$ is a module for
$H_{G,T}^{\mathrm{lv}}$.  This follows directly from \eqref{eq:L1},
using $g^{-1}H_{G,T}^{\mathrm{lv}}g\subseteq H_{G,T}^0$.
The convolution uses the positive-loop descended kernel with $!$
pairing and star integration; its unmodified-section value is $M$.

The level framing defines the modification map
\[
 t_Z:Z\times_S B_G^{\mathrm{lv}}\longrightarrow Y_G\times S.
\]
The determinant formula for this map has the factor $\mathcal L^s$
on $Z$, which cancels the twist of $K$.

\begin{proposition}\label{loc:kernel}
There is a natural isomorphism
\begin{equation}\label{eq:T6}
 \operatorname{Loc}_{G,T}^{\mathrm{lv}}(K\star M)
 \simeq
 (\operatorname{pr}_{B_G^{\mathrm{lv}}})_*
 \left(\operatorname{pr}_Z^!K\otimes^!
                         t_Z^!\operatorname{Loc}_T(M)\right).
\end{equation}
It is compatible with ordinary and crystalline parameter operations,
with deeper levels, and with insertion into $E\supseteq T$.  In the
last comparison, $K$ is taken by closed direct image to the
Grassmannian at $E$.
\end{proposition}

\begin{proof}
\emph{1. Relative gauge comparison.}
We first compare the kernels before integrating against $K$.  On a
smooth affine plot of $Z$, choose a Laurent lift $g$ and a determinant
lift $\widetilde g$.  If $f$ is a disk frame on a localization chart,
the frames $f$ and $fg$ identify their vector/Lie pairs by
\begin{equation}\label{loc:gauge}
 (v,x)\longmapsto
 \left(v,\operatorname{Ad}_{g^{-1}}x+
                   \widetilde g^{-1}v(\widetilde g)\right).
\end{equation}
This is the change in \eqref{eq:L1}, including the determinant
correction.  Hence \eqref{eq:T2} identifies the relative kernel
formulas at fixed plot parameters.  To obtain a comparison of
D-module kernels, we must also verify formal transport when the plot
parameter moves.

\emph{2. Crystalline transfer and independence of choices.}
Consider two parameter lifts across a fixed nilpotent ideal.  Use
common full frames, passing to a common deeper level if necessary.
After smooth chart refinement, the two modified bundles lift to maps
$q_1,q_2:W\to Y$ to the same smooth bundle chart, agreeing on reduction.
Choose their disk and determinant frames with matching reductions;
other choices are handled by the arc descent already proved.  The
change between the two lifts is a formal arrow of the replacement
groupoid of $A_Y=\mathcal T_Y$.  It consists of the nearby chart maps,
a trivialization of the deformation off the markings, and its Laurent
replacement in the chosen disk frames, including the determinant lift.

We describe its action to all nilpotent orders.  Locally split the
transitive anchor of $A_Y$.  Lifts of coordinate derivatives together
with outside-current directions give coordinates for formal arrows:
first move the target along the lifted coordinate directions, then
apply a formal outside automorphism.  The first movement may be
constructed recursively by its formal flow, integrating the Laurent
vector and its central component through the logarithmic derivative.
It remains a replacement arrow because the framed gluing relation is
preserved at each nilpotent order.  Equivalently, construct the outside
trivializations order by order; their obstruction vanishes by the
affineness of $X-E$.  Formal isotropy is the exponential of outside
currents, with the indicated central splitting.

The algebra of finite tangent distributions on this formal groupoid
is therefore $U(A_Y)$.  Indeed its order filtration has associated
graded the symmetric powers of the coordinate and outside-current
tangents, giving the usual Taylor proof of PBW.  Only finite Taylor
expressions are needed: the coordinate displacements are nilpotent,
and at each fixed nilpotent order a Laurent section has a finite pole
bound and uses finitely many tangent expressions.  No dual of all
functions on an infinite-dimensional formal germ is used.  Formal
composition gives the composition of these distributions.

For clarity, this identifies the derived transfer maps as well.
The two module pulls use
\[
 \mathcal O_W\otimes_{\mathcal O_Y,q_i}U(A_Y),\qquad i=1,2,
\]
as right $U(A_Y)$-modules with their inverse-image algebroid actions.
Multiply by the Taylor distribution of the formal arrow to send the
second generating section to the first transfer.  This map is right
linear.  Differentiating composition shows that it intertwines the
inverse-image algebroids by \eqref{loc:gauge}.  After tensoring with
the chart TDO, the same distribution becomes the ordinary TDO Taylor
transfer, with the determinant-frame factor.  A nil-unit change of
that lift acts by its inverse-section power in the arrow from the
second pull to the first, on both sides, because the central scalar is
$-s$.

One can verify the transfer quotient directly by PBW.  The anchor
kernel pulls back from $Y$; order its factors first.  The enveloping
transfer is then PBW-free in these kernel directions, so tensoring to
the kernel quotient computes its derived tensor as well.  Equivalently,
the associated-graded Koszul complex has no additional homology.
The resulting quotient is precisely the ordinary TDO transfer.

The result is independent of the outside trivialization used to
construct the formal arrow.  Two choices differ by a formal outside
automorphism.  On a fixed nilpotent test its Taylor distribution is
a finite unit $u$ in the enveloping algebra of the target
inverse-image algebroid, lying in the anchor-kernel subalgebra.
Its determinant-trivial lift maps to $1$ in the TDO quotient.
The inner-conjugation bimodule is canonically identified by
multiplication by $u$, with $u^{-1}$ for the reverse direction;
this intertwines the corresponding coefficient-module actions.
On the TDO quotient the identification is the identity.  These
bimodule identifications compose by multiplication of units.
Applying the functorial two-sided bar construction therefore gives
a coherent trivialization of the ambiguity after derived tensor,
including all higher homotopies.

Thus tensoring these transfer diagrams with a complex of coefficient
modules proves the required comparison, including derived relations
in compact presentations.  On the coefficient module the formal
arrow acts by the Taylor form of \eqref{eq:L1}, since its derivatives
are the stated Lie action.  Formal-arrow composition gives every
cocycle identity.  Refining charts or changing arc frames uses the
same transfer and Harish--Chandra formulas.

\emph{3. Integration and insertion.}
We have proved equality after every $!$-plot pull, with all the
infinitesimal lift comparisons.  The plot criterion proved in
Proposition~\ref{loop:strong-action} therefore identifies the parameter D-module kernels.  Localization
with parameter coefficients is continuous and $D(S)$-linear.
Star-pairing these kernels with $K$ commutes with localization, by
relative tensor and $!$ base change, and gives \eqref{eq:T6}.  The
calculation is schematic on finite bounds, so it is valid on the
smooth bundle charts used above.

Finally insert into $E$.  The modification represented by $g$ has
poles only at $T$.  In an $E$-frame its Laurent datum consequently
belongs to the intermediate meromorphic group: clear powers of the
old divisor before completing, as in
Proposition~\ref{loc:propagation}.  Conjugation identifies the
intermediate-algebra induction bars with the modified old-module
bars.  Termwise their Lie actions are intertwined by
\eqref{loc:gauge}, with the same relative determinant.  The derivative
on the added enveloping factor is the conjugation derivative; together
with the derivative on the coefficient module this is exactly
\eqref{eq:L1} on the induced module.  Hence the preceding formal Taylor
comparison also intertwines the bars.  The old level depth suffices:
only poles at the old divisor are cleared, and the new positions
carry positive arcs.  Arc changes use Harish--Chandra equivariance.
The plot comparison and star integration now prove the insertion
assertion.  At any enlarged finite bound, imposing the old pole
conditions is a closed embedding, so closed-image/$!$-pull formulas
supply the stated image of $K$.  All these comparisons are unchanged
at deeper levels.
\end{proof}

\subsection{Whittaker averaging as a unipotent period}

Localization and its kernel compatibility are now available in
families.  We use them to express the left side of \eqref{eq:T1},
before the scalar normalization in $\Lambda_T$, as one exponential
period.

Let $q_N^{\mathrm{lv}}:B_N^{\mathrm{lv}}\to S$ be projection and let
$\psi_0$ be the zero-defect function in \eqref{eq:W1}, pulled back to
this stack.  Define
\[
 \begin{split}
 \operatorname{Per}_{G,E}^{\mathrm{lv}}(M)
   &=(q_N^{\mathrm{lv}})_*
     \left(\operatorname{Exp}(-\psi_0)\otimes^!
                  p^!\operatorname{Loc}_{G,E}^{\mathrm{lv}}(M)\right),\\
 \operatorname{Per}_{N,E}^{\mathrm{lv}}(M)
   &=(q_N^{\mathrm{lv}})_*
     \left(\operatorname{Exp}(-\psi_0)\otimes^!
                    \operatorname{Loc}_{N,E}^{\mathrm{lv}}(M)\right).
 \end{split}
\]
Their projections are safe, since stabilizers are unipotent, and these
functors have the preceding parameter compatibilities.  Put
\[
 \delta=\dim(B_N^{\mathrm{lv}}/S)
             =-\chi(X,\mathfrak n^{\mathrm{lv}}).
\]
The periods with shift $[-2\delta]$ agree under deeper nonzero-level
pull comparisons.  This is the smooth normalization for the
unipotent torsors that forget the extra level.

For $W'=\mathbf d(K)$, where $K$ is supported on $Z$, the pairing
\eqref{eq:loop-whittaker-functional} expresses the unnormalized left side
as
\begin{equation}\label{eq:T7}
 (q_Z)_*\left(K\otimes^!\mu^!R_T
                              \operatorname{Loc}_T(M)\right).
\end{equation}
The maps are $q_Z:Z\to S$ and the uniformization map $\mu$ used in
the Whittaker model.

\begin{proposition}\label{loc:averaging}
For the deep level chosen in Proposition~\ref{loc:kernel},
Expression~\eqref{eq:T7} is naturally isomorphic to
\begin{equation}\label{eq:T8}
 \operatorname{Per}_{G,T}^{\mathrm{lv}}(K\star M)[-2\delta].
\end{equation}
The isomorphism respects parameter operations, increases of finite
bounds, and deeper-level comparisons.
\end{proposition}

\begin{proof}
We first recover $\mu^!R_T(F)$ on $Z$, for any family $F$ of bundle
D-modules, by finite-loop averaging.  We then pair this formula with
$K$ and exchange the two integrations; the dimensions of the unipotent
presentation give the period shift.

Let $L_l\subset LN_T$ be the compact subgroup
allowing pole order $l\operatorname{ht}(\alpha)$ times the marking
divisor in each root direction $\alpha$.  The divisor counts repeated
marks.  For fixed $l$ and $Z$, a sufficiently deep normal congruence
$J\subset L_l$ makes the action, determinant, and character formulas
factor through a finite-dimensional quotient.  Normalized star
averaging evaluated at $g$ integrates the bundle pull at $ng$
against $\operatorname{Exp}(-\chi(n))$ over that quotient, with shift
minus twice its relative dimension.  Passing to a deeper quotient
leaves the result unchanged by unipotent de Rham acyclicity.  The
$N$-equivariance of $R_T(F)$ and its counit to the bundle pull give a
natural comparison from $\mu^!R_T(F)$ to this averaging expression.

We show that the comparison is an isomorphism for sufficiently large
$l$.  Stratify $Z$ first by coincidences and then by semi-infinite orbit
type.  Only finitely many types occur: the Pl\"ucker orders are bounded
above and below by the lattice bounds in a finite collection of
representations.  On each stratum, divide the Pl\"ucker maps by their
prescribed zero orders.  This produces the genuine Borel extension
with its Cartan lattices.  Choose unipotent frames on the disks and
compare to the punctured trivialization.  Locally on the parameter
stratum the result has the form $n_0t^\lambda$, with $\lambda$ not
assumed dominant.  These representatives have bounded Laurent poles;
a finite cover of the bounded strata gives a common allowance $l$.
The calculation ignores nilpotent structure only when forming the
stratification; the comparison of D-modules is pulled by $!$ base
change, so it applies with arbitrary parameter coefficients.

The stabilizer of $t^\lambda$ is
\[
 J_\lambda=t^\lambda H_{N,T}^0t^{-\lambda}\subseteq L_l
\]
for sufficiently large $l$.  Here the superscript $0$ means zero level
divisor, hence full arcs.  The residual integration along its finite
unipotent quotient is zero unless $\chi|_{J_\lambda}=0$: the de Rham
cohomology of a nontrivial exponential on an additive root quotient
vanishes.  Testing the simple roots shows that this condition is
exactly dominance of $\lambda$.  The determinant on the stabilizer
has its canonical splitting.

For dominant $\lambda$, the quotient $L_l/J_\lambda$ presents the
genuine bundle data with Cartan bundle
$T_\lambda=T_0(-\sum_y\lambda_y y)$, modulo a smooth finite-dimensional
unipotent group.  To see this in families, enlarge the root lattices
at the markings to those defining $L_l$, leaving the group unchanged
on their complement.  On each locally closed parameter stratum $S_D$,
choose $l$ so that
$R^1\pi_*$ vanishes for every enlarged root line, where
$\pi:X\times S_D\to S_D$.  Then their $\pi_*$ are locally free and
commute with base change.  Successive vector extensions show that
the enlarged group bundles are locally trivial in families; their trivializations form a torsor for the
unipotent group of global sections.  Recovering the original lattice
is exactly the quotient $L_l/J_\lambda$.  The residue convention
makes the exponential descend to this presentation.

Enlarge $l$ once more so that the chosen $n_0$ lies in $L_l$.
Pulling this presentation to the stratum identifies normalized
averaging with the conjugated pull after star image along the
stratum map $q_D$ in \eqref{eq:W1}.  Indeed, $!$ base change gives
that correspondence; integration along its smooth unipotent
presentation fibers, together with the averaging normalization,
gives precisely $[-2\dim_{\mathrm{rel}}q_D]$.  The averaging correspondence
and the Whittaker stratum projector
in \eqref{eq:W1} are thus $!$-base changes of the same relative quotient
presentation.  Their counits induce the identity of that projector,
with the displayed dimension shift.  This proves the comparison for
arbitrary coefficient D-modules on $S_D$.  The strata are finite, so
their $!$ restrictions
prove the assertion on $Z$.  Increasing $l$ uses the same averaging
projection and trace, hence gives the same map.

Apply this calculation inside \eqref{eq:T7}.  The level choice allows
us to project the averaging variable to
$L_l/H_{N,T}^{\mathrm{lv}}$.  Its element represents a point of
$B_N^{\mathrm{lv}}$, and the product $ng$ is precisely the level
modification in \eqref{eq:T6}.  The line and phase descend with this
map.  Integration over the unipotent difference changes the
normalization to $[-2e_l]$, where
\[
 e_l=\dim(L_l/H_{N,T}^{\mathrm{lv}}).
\]
By the preceding enlarged-root argument, this quotient presents
$B_N^{\mathrm{lv}}$ modulo the unipotent group of global sections of
the enlarged bundle.  If that group has dimension $f_l$, then
$\delta=e_l-f_l$.  Integration over its presentation fibers changes
$[-2e_l]$ to $[-2e_l+2f_l]=[-2\delta]$.
The descended character is $\operatorname{Exp}(-\psi_0)$, and the
determinant splittings agree by multiplicativity along the root
filtration.  The remaining integration against $K$ is exactly
\eqref{eq:T6}, proving \eqref{eq:T8}.

All shifts here follow the fixed smooth/left convention: integrating
a dualizing pull over an affine-space fiber of dimension $e$
contributes $[2e]$.  Every calculation takes place at finite bounds
and finite congruence quotients, and the comparison maps are the
normalized pull, trace, and base-change maps.  Their compatibility
under refinement proves the asserted naturality.
\end{proof}
